\documentclass[11pt, a4paper]{article}

% ============================================================
% Pacchetti fondamentali
% ============================================================
\usepackage[utf8]{inputenc}
\usepackage[italian]{babel}
\usepackage[T1]{fontenc}
\usepackage{lmodern} % Ottimizza la resa dei font in PDF
\usepackage{amsmath} % Comandi matematici (\text, \rightarrow, ...)

% Formattazione pagina e layout
\usepackage{geometry}
\geometry{top=2.5cm, bottom=2.5cm, left=2.5cm, right=2.5cm}

% Tabelle professionali
\usepackage{tabularx}
\usepackage{booktabs}
\usepackage{array}

% Grafica, colori, elenchi, codice
\usepackage{graphicx}
\usepackage{xcolor}
\usepackage{enumitem}
\usepackage{listings}

% Intestazioni e piè di pagina
\usepackage{fancyhdr}
\pagestyle{fancy}
\fancyhf{}
\fancyhead[L]{\textcolor{gray}{\textbf{Travel Organizer Offline}}}
\fancyhead[R]{\textcolor{gray}{Relazione Tecnica}}
\fancyfoot[C]{\thepage}
\renewcommand{\headrulewidth}{0.4pt}
\renewcommand{\footrulewidth}{0.4pt}
\setlength{\headheight}{14pt} % Evita il warning di fancyhdr sull'altezza

% Stile dei titoli
\usepackage{titlesec}
\titleformat{\section}{\Large\bfseries\color{blue!70!black}}{\thesection}{1em}{}
\titleformat{\subsection}{\large\bfseries\color{blue!50!black}}{\thesubsection}{1em}{}
\titleformat{\subsubsection}{\normalsize\bfseries\color{blue!30!black}}{\thesubsubsection}{1em}{}

\usepackage[colorlinks=true, linkcolor=blue!70!black, urlcolor=blue!70!black]{hyperref}

% Stile per i frammenti di codice
\definecolor{codebg}{RGB}{245,247,250}
\lstset{
  basicstyle=\ttfamily\footnotesize,
  backgroundcolor=\color{codebg},
  frame=single,
  rulecolor=\color{gray!40},
  breaklines=true,
  columns=fullflexible,
  keepspaces=true,
  showstringspaces=false
}

\begin{document}

% ============================================================
% Frontespizio
% ============================================================
\begin{titlepage}
    \centering
    \vspace*{2.5cm}
    {\Huge \textbf{Relazione Tecnica}}\\[0.4cm]
    {\Large di Progettazione e Sviluppo}\\[1.5cm]
    {\Huge \textbf{Travel Organizer Offline}}\\[0.5cm]
    {\large App mobile per l'organizzazione dei viaggi}\\[2cm]

    {\Large Giugno 2026\par}
    \vfill
    \begin{flushright}
        \textbf{\textit{Gruppo}}\\
        \textit{Componente 1}\\
        \textit{Componente 2}\\
        \textit{Componente 3}\par
    \end{flushright}
    \vspace{2cm}
    \Large Mobile Programming -- Anno Accademico 2025/2026
\end{titlepage}

\tableofcontents
\newpage

% ============================================================
\section{Descrizione dell'app}

\subsection{Obiettivo}
\textbf{Travel Organizer Offline} è un'applicazione mobile, sviluppata in
\textit{Flutter}, che aiuta l'utente a pianificare e gestire i propri viaggi:
itinerari, tappe, attività, checklist e spese. Tutte le informazioni restano
disponibili \textbf{anche senza connessione di rete}, perché i dati sono salvati
in un database locale sul dispositivo. L'app non utilizza API remote né un
backend: è pensata per funzionare al 100\% offline.

\subsection{Tipologia di utenti}
L'applicazione si rivolge a chiunque organizzi viaggi -- viaggiatori abituali,
studenti, famiglie -- e desideri uno strumento semplice per tenere sotto
controllo programma, bagaglio e budget anche dove la rete non è disponibile
(aereo, estero, zone remote).

\subsection{Problema che l'app intende risolvere}
Le informazioni di un viaggio sono spesso frammentate tra note, chat e siti
diversi e richiedono una connessione per essere consultate. Travel Organizer
Offline centralizza tutto in un'unica app consultabile offline, evitando la
dipendenza dalla rete e la dispersione dei dati.

\subsection{Principali scenari d'uso}
\begin{itemize}[noitemsep]
    \item \textbf{Pianificazione:} l'utente crea un viaggio, lo suddivide in
    tappe e associa a ciascuna le attività con orari e luoghi.
    \item \textbf{Preparazione:} genera la lista valigia in base ai tag del
    viaggio e compila le checklist dei documenti.
    \item \textbf{Controllo spese:} registra spese previste ed effettive e
    monitora il budget tramite la barra di salute.
    \item \textbf{In viaggio:} consulta offline la timeline della giornata e i
    tempi liberi disponibili.
    \item \textbf{Condivisione:} esporta un viaggio in JSON e lo invia a un
    amico, che può reimportarlo nella propria app.
\end{itemize}

% ============================================================
\section{Requisiti}

\subsection{Quadro generale dei requisiti}
\begin{table}[h!]
    \centering
    \renewcommand{\arraystretch}{1.3}
    \begin{tabularx}{\textwidth}{>{\raggedright\arraybackslash}p{3cm} l X l}
        \toprule
        \textbf{Categoria} & \textbf{Codice} & \textbf{Descrizione} & \textbf{Priorità} \\
        \midrule
        Dati & DF-1 & Modellazione e persistenza (Viaggi, Tappe, Attività, Checklist, Spese) & Must \\
        Funzione & IF-1 & Gestione dei Viaggi & Must \\
        Funzione & IF-2 & Gestione delle Tappe/Giornate & Must \\
        Funzione & IF-3 & Gestione delle Attività e dell'itinerario & Must \\
        Funzione & IF-4 & Gestione delle Checklist e informazioni utili & Must \\
        Funzione & IF-5 & Gestione delle Spese e del budget & Must \\
        Funzione & IF-6 & Ricerca, filtri e organizzazione & Should \\
        Funzione & IF-7 & Riepiloghi e statistiche & Should \\
        Funzione & IF-8 & Feature avanzate & Should \\
        Flusso & BF & Flussi inter-modulo & Must \\
        \bottomrule
    \end{tabularx}
\end{table}

\subsection{Modelli dei dati (DF)}
\begin{itemize}[noitemsep]
    \item \textbf{DF-1.1 Viaggio:} id, titolo, destinazione, dataInizio,
    dataFine, descrizione, stato, budget, partecipanti, note, tag, dataCreazione.
    \item \textbf{DF-1.2 Tappa:} id, viaggioId, titolo, data, località,
    descrizione, ordine, note.
    \item \textbf{DF-1.3 Attività:} id, viaggioId, tappaId, titolo, descrizione,
    dataOra, luogo, categoria, costoPrevisto, stato, note.
    \item \textbf{DF-1.4 Checklist:} id, viaggioId, tappaId, titolo, descrizione,
    elementi.
    \item \textbf{DF-1.5 Elemento checklist:} id, checklistId, titolo,
    completato, ordine.
    \item \textbf{DF-1.6 Spesa:} id, viaggioId, tappaId, attivitàId, titolo,
    importo, categoria, data, metodoPagamento, stato (prevista/effettiva), note.
\end{itemize}

\subsection{Funzionalità implementate (IF)}
\begin{itemize}[noitemsep]
    \item \textbf{IF-1 Viaggi:} elenco, dettaglio, inserimento, modifica,
    eliminazione, \textbf{duplicazione} e assegnazione di \textbf{tag}. Lo stato
    (futuro/in corso/completato/archiviato) è calcolato dalle date.
    \item \textbf{IF-2 Tappe:} elenco, inserimento, modifica, eliminazione,
    \textbf{riordino con drag \& drop} e accesso alle attività associate.
    \item \textbf{IF-3 Attività:} inserimento, modifica, eliminazione,
    associazione a viaggio/tappa, cambio di stato (da fare/completata/annullata),
    categoria, data/ora, luogo e costo previsto.
    \item \textbf{IF-4 Checklist:} creazione, aggiunta/modifica/eliminazione e
    spunta degli elementi, barra di avanzamento e duplicazione.
    \item \textbf{IF-5 Spese:} inserimento, modifica, eliminazione, distinzione
    tra spese previste ed effettive, metodo di pagamento e confronto col budget.
    \item \textbf{IF-6 Ricerca e filtri:} viaggi per titolo/destinazione e per
    stato; attività per categoria, stato e giorno; spese per stato, categoria e
    fascia di importo; tappe per nome/località e per data.
    \item \textbf{IF-7 Statistiche:} totali e stato dei viaggi, attività totali e
    completate, spese effettive/previste, avanzamento checklist, classifica delle
    tappe con più attività, grafico a torta delle spese per categoria e barra di
    salute del budget per ciascun viaggio.
    \item \textbf{IF-8 Feature avanzate:} descritte nella Sezione~\ref{sec:adv}.
\end{itemize}

\subsection{Feature avanzate scelte}\label{sec:adv}
La traccia richiede almeno due feature avanzate; il progetto ne integra diverse,
tutte effettivamente funzionanti:
\begin{enumerate}[noitemsep]
    \item \textbf{Packing list intelligente con tag.} A ogni viaggio si possono
    assegnare dei tag (Mare, Montagna, Città, Estero, Business, Trekking,
    Inverno). Una matrice predefinita nel codice associa a ogni tag un insieme di
    oggetti (es. \textit{Mare} $\rightarrow$ Costume, Crema solare; \textit{Estero}
    $\rightarrow$ Passaporto, Adattatore prese). La generazione della lista valigia
    unisce gli oggetti di base con i suggerimenti dei tag, senza duplicati. Tutta
    la logica è locale.
    \item \textbf{Timeline visuale con tempi liberi.} L'itinerario è mostrato come
    linea del tempo con blocchi colorati in base alla categoria; l'app calcola i
    periodi liberi tra un'attività e l'altra e li evidenzia (es. \textit{``Hai 3h di
    tempo libero''}).
    \item \textbf{Budget gamification.} Una barra di salute del budget cambia
    colore (verde $\rightarrow$ giallo $\rightarrow$ rosso) all'aumentare della
    spesa, e un avviso compare quando una spesa effettiva supera il budget residuo.
    \item \textbf{Esporta/Importa viaggio.} Un viaggio completo (con tappe,
    attività, checklist e spese) può essere esportato in JSON e condiviso tramite
    il menu nativo del sistema; un altro utente può reimportarlo. Gli identificatori
    vengono rigenerati per evitare conflitti con i dati esistenti.
    \item \textbf{Duplicazione} di viaggi, tappe e checklist.
\end{enumerate}

\subsection{Funzionalità considerate ma non implementate}
\begin{itemize}[noitemsep]
    \item Simulatore di cambio valuta con tasso manuale per viaggio.
    \item Salvataggio di ``giornate tipo'' come modelli riutilizzabili.
    \item Temi grafici dinamici in base alla destinazione.
    \item Notifiche locali (promemoria di partenza/scadenze).
\end{itemize}
Queste funzioni sono state valutate ma escluse per mantenere l'app semplice,
coerente e priva di dipendenze non necessarie.

\subsection{Limitazioni note}
\begin{itemize}[noitemsep]
    \item La gestione del cambio valuta non è inclusa: gli importi sono in un'unica
    valuta (Euro).
    \item L'import di un viaggio avviene incollando il testo JSON (non è prevista la
    lettura diretta di un file).
    \item I dati sono locali al dispositivo: non esiste sincronizzazione su cloud.
\end{itemize}

% ============================================================
\section{Progettazione dell'app}

\subsection{Struttura generale}
L'app adotta un'architettura a livelli che separa nettamente
\textbf{interfaccia}, \textbf{stato} e \textbf{accesso ai dati}:
\[
\text{UI (Schermate)} \;\rightarrow\; \text{Provider (Stato)} \;\rightarrow\;
\text{Repository} \;\rightarrow\; \text{Database SQLite}
\]
Questa separazione rende il codice ordinato e testabile: le schermate non
conoscono il database, ma dialogano con i provider; i provider usano i repository,
che incapsulano le query SQL.

\subsection{Schermate principali}
\begin{itemize}[noitemsep]
    \item \textbf{Lista Viaggi:} elenco con ricerca, filtro per stato, tag e azioni
    rapide (apri, duplica, modifica, elimina); pulsante di importazione.
    \item \textbf{Dettaglio Viaggio:} sei schede -- Tappe, Attività, Checklist,
    Spese, Timeline, Valigia -- con esportazione/condivisione.
    \item \textbf{Form} di inserimento/modifica per viaggi, tappe, attività e spese,
    con validazione dei campi.
    \item \textbf{Dettaglio Tappa:} informazioni della tappa ed elenco attività.
    \item \textbf{Statistiche:} riepiloghi e grafici su tutti i viaggi.
\end{itemize}

\subsection{Flusso di navigazione}
La navigazione principale avviene tramite una \textit{Bottom Navigation Bar} a due
sezioni (Viaggi e Statistiche). Dalla lista si apre il dettaglio di un viaggio
(passaggio dell'\texttt{id} alla schermata successiva), organizzato in schede; da
ciascuna scheda si raggiungono i form di inserimento/modifica. Le transizioni usano
\texttt{Navigator} e \texttt{MaterialPageRoute}, con passaggio esplicito dei dati
tra schermate.

\subsection{Organizzazione dei dati}
Lo schema relazionale prevede una tabella per ciascuna entità, con vincoli di
chiave esterna e \textbf{eliminazione a cascata}: cancellando un viaggio vengono
rimossi automaticamente tappe, attività, checklist (con i relativi elementi) e
spese collegate.

\begin{table}[h!]
    \centering
    \renewcommand{\arraystretch}{1.3}
    \begin{tabularx}{\textwidth}{l l X}
        \toprule
        \textbf{Tabella} & \textbf{Relazione} & \textbf{Note} \\
        \midrule
        trips & -- & Entità principale \\
        stages & trips (1:N) & ON DELETE CASCADE \\
        activities & trips/stages & CASCADE su trip, SET NULL su stage \\
        checklists & trips/stages & CASCADE su trip \\
        checklist\_items & checklists (1:N) & ON DELETE CASCADE \\
        expenses & trips (1:N) & ON DELETE CASCADE \\
        \bottomrule
    \end{tabularx}
\end{table}

\subsection{Integrità dei dati}
Oltre ai vincoli del database, l'app applica regole a livello applicativo:
\begin{itemize}[noitemsep]
    \item impedisce di salvare due viaggi con \textbf{date sovrapposte};
    \item valida che la data di fine non preceda quella di inizio;
    \item richiede i campi obbligatori (titolo, destinazione, importo);
    \item dopo l'eliminazione di un viaggio, allinea le cache in memoria così le
    statistiche restano coerenti.
\end{itemize}

% ============================================================
\section{Scelte tecnologiche}

\subsection{Framework}
È stato scelto \textbf{Flutter} (Material 3) per la possibilità di sviluppare con
un unico codice sorgente un'app fluida e multipiattaforma (Android, iOS, Web),
coerentemente con i vincoli della traccia.

\subsection{Librerie principali}
\begin{table}[h!]
    \centering
    \renewcommand{\arraystretch}{1.3}
    \begin{tabularx}{\textwidth}{l X}
        \toprule
        \textbf{Libreria} & \textbf{Motivazione} \\
        \midrule
        \texttt{provider} & Gestione dello stato semplice e reattiva (pattern ChangeNotifier). \\
        \texttt{sqflite} / \texttt{sqflite\_common\_ffi\_web} & Persistenza locale su SQLite, anche su Web tramite WebAssembly. \\
        \texttt{fl\_chart} & Grafici (torta) per la sezione statistiche. \\
        \texttt{intl} & Formattazione di date e valuta in italiano. \\
        \texttt{share\_plus} & Condivisione tramite il menu nativo del sistema. \\
        \texttt{uuid} & Generazione di identificatori univoci. \\
        \bottomrule
    \end{tabularx}
\end{table}

\subsection{Motivazione delle scelte}
\texttt{provider} è leggero e adatto a un'app di queste dimensioni, evitando la
complessità di soluzioni più pesanti. \texttt{sqflite} fornisce un database
relazionale completo (chiavi esterne, cascata) ideale per dati strutturati e
collegati. La combinazione con \texttt{sqflite\_common\_ffi\_web} permette di
eseguire l'app anche nel browser senza modifiche al codice di dominio.

\subsection{Gestione della persistenza}
Tutti i dati sono salvati in un database \textbf{SQLite locale}. Un
\texttt{DatabaseHelper} (singleton) gestisce apertura, creazione dello schema,
attivazione delle chiavi esterne e migrazioni di versione. I \textbf{Repository}
incapsulano le query SQL (CRUD e ricerche) per ciascuna entità.

% ============================================================
\section{Implementazione}

\subsection{Organizzazione del codice}
\begin{lstlisting}
lib/
  core/        Tema, colori, utility, catalogo packing
  data/
    database/  DatabaseHelper (schema, migrazioni)
    models/    Trip, Stage, Activity, Checklist, Expense
    repositories/  Accesso ai dati (query SQL)
    trip_transfer.dart  Esporta/Importa viaggio (JSON)
  providers/   Stato dell'app (ChangeNotifier)
  features/    Schermate (viaggi, tappe, attivita', spese, timeline, statistiche)
  shared/widgets/  Widget riutilizzabili
  main.dart    Avvio dell'app
\end{lstlisting}

\subsection{Widget principali}
Tra i componenti riutilizzabili: \texttt{EmptyState} (schermata per le liste
vuote con call-to-action), \texttt{ConfirmDialog} (conferma delle eliminazioni),
\texttt{StatusChip} (etichetta colorata di stato). Le schermate di dettaglio usano
\texttt{TabBar}/\texttt{TabBarView} e \texttt{SliverAppBar}.

\subsection{Gestione dello stato}
Lo stato è gestito con il pattern \textbf{Provider/ChangeNotifier}: ogni entità ha
un proprio provider (es. \texttt{TripProvider}, \texttt{ExpenseProvider}) che
mantiene i dati in memoria, esegue le operazioni tramite i repository e notifica la
UI con \texttt{notifyListeners()}. Le schermate si aggiornano automaticamente
tramite \texttt{Consumer} e \texttt{context.watch}.

\subsection{Gestione della navigazione}
La navigazione tra schermate usa \texttt{Navigator} e \texttt{MaterialPageRoute},
con \textbf{passaggio esplicito dei dati} (ad esempio l'\texttt{id} del viaggio o
l'oggetto da modificare). La barra inferiore usa un \texttt{IndexedStack} per
conservare lo stato delle sezioni.

\subsection{Gestione dei dati persistenti}
Le operazioni di lettura/scrittura passano sempre dai repository. Esempio di
controllo di integrità sulle date dei viaggi:
\begin{lstlisting}
// Due intervalli si sovrappongono se: inizio <= altraFine && altroInizio <= fine
Trip? overlappingTrip(DateTime start, DateTime end, {String? excludeId}) {
  for (final t in _trips) {
    if (t.id == excludeId) continue;
    if (!start.isAfter(t.endDate) && !t.startDate.isAfter(end)) return t;
  }
  return null;
}
\end{lstlisting}

\subsection{Parti significative}
\begin{itemize}[noitemsep]
    \item \textbf{Packing list intelligente:} un catalogo \texttt{tag $\rightarrow$
    oggetti} guida la generazione automatica della valigia.
    \item \textbf{Esporta/Importa:} serializzazione dell'intero grafo di un viaggio
    in JSON e re-importazione con rimappatura degli identificatori, validata da un
    test automatico di round-trip.
    \item \textbf{Timeline:} calcolo dei tempi liberi tra attività consecutive dello
    stesso giorno.
    \item \textbf{Test:} il progetto include test automatici (avvio dell'app,
    integrità delle date, eliminazione, esporta/importa) eseguiti su SQLite via FFI.
\end{itemize}

% ============================================================
\section{Casi d'uso}

\subsection{Elenco dei casi d'uso}
\begin{table}[h!]
    \centering
    \renewcommand{\arraystretch}{1.3}
    \begin{tabularx}{\textwidth}{l X | l X}
        \toprule
        \textbf{ID} & \textbf{Caso d'uso} & \textbf{ID} & \textbf{Caso d'uso} \\
        \midrule
        UC-1 & Inserisci viaggio & UC-7 & Cerca e filtra elementi \\
        UC-2 & Inserisci tappa & UC-8 & Visualizza statistiche \\
        UC-3 & Pianifica attività & UC-9 & Visualizza timeline \\
        UC-4 & Gestisci checklist & UC-10 & Esporta / Importa viaggio \\
        UC-5 & Registra spesa & UC-11 & Modifica o elimina elemento \\
        UC-6 & Genera packing list smart & & \\
        \bottomrule
    \end{tabularx}
\end{table}

\subsection{Descrizione dettagliata}

\subsubsection*{UC-1: Inserisci viaggio}
\textbf{Attore:} Utente.\\
\textbf{Precondizioni:} App avviata, persistenza inizializzata.\\
\textbf{Postcondizioni:} Un nuovo viaggio è salvato e mostrato nell'elenco.\\
\textbf{Flusso base:}
\begin{enumerate}[noitemsep]
    \item L'utente seleziona ``Nuovo viaggio''.
    \item Compila titolo, destinazione, date, budget, partecipanti, tag e note.
    \item Il sistema valida i campi obbligatori e le date.
    \item Il sistema salva il viaggio e torna all'elenco aggiornato.
\end{enumerate}
\textbf{Flussi alternativi:}
\begin{itemize}[noitemsep]
    \item \textbf{3a. Date sovrapposte:} il sistema rileva un altro viaggio con date
    che si accavallano, blocca il salvataggio e mostra un avviso.
    \item \textbf{3b. Campi non validi:} il sistema evidenzia l'errore e resta in
    attesa di una correzione.
\end{itemize}

\subsubsection*{UC-2: Inserisci tappa}
\textbf{Attore:} Utente.\\
\textbf{Precondizioni:} Esiste un viaggio.\\
\textbf{Postcondizioni:} La tappa è aggiunta all'itinerario del viaggio.\\
\textbf{Flusso base:} l'utente apre la scheda Tappe, inserisce titolo, data,
località, descrizione e note; il sistema valida e salva, aggiornando l'elenco
(riordinabile tramite drag \& drop).

\subsubsection*{UC-3: Pianifica attività}
\textbf{Attore:} Utente.\\
\textbf{Precondizioni:} Esiste un viaggio (ed eventualmente una tappa).\\
\textbf{Postcondizioni:} L'attività è pianificata e associata al viaggio/tappa.\\
\textbf{Flusso base:} l'utente inserisce titolo, categoria, data/ora, luogo, costo
e tappa associata; il sistema valida e salva. \textbf{Alternativo:} titolo mancante
$\rightarrow$ il salvataggio è bloccato con messaggio d'errore.

\subsubsection*{UC-4: Gestisci checklist}
\textbf{Attore:} Utente.\\
\textbf{Postcondizioni:} La checklist e i suoi elementi sono aggiornati.\\
\textbf{Flusso base:} l'utente crea una checklist, aggiunge elementi, li spunta o
li elimina; la barra di avanzamento si aggiorna in tempo reale. Può inoltre
duplicare una checklist esistente.

\subsubsection*{UC-5: Registra spesa}
\textbf{Attore:} Utente.\\
\textbf{Postcondizioni:} La spesa è registrata e i totali aggiornati.\\
\textbf{Flusso base:} l'utente inserisce titolo, importo, categoria, data, tipo
(prevista/effettiva) e metodo di pagamento; il sistema valida e salva.\\
\textbf{Alternativo 3a. Fuori budget:} se la spesa effettiva supera il budget
residuo, il sistema mostra un avviso e chiede conferma prima di salvare.

\subsubsection*{UC-6: Genera packing list smart}
\textbf{Attore:} Utente.\\
\textbf{Precondizioni:} Esiste un viaggio (preferibilmente con tag).\\
\textbf{Postcondizioni:} La lista valigia è popolata con oggetti pertinenti.\\
\textbf{Flusso base:}
\begin{enumerate}[noitemsep]
    \item L'utente apre la scheda Valigia e seleziona ``Genera lista''.
    \item Il sistema legge i tag del viaggio e recupera dal catalogo gli oggetti
    suggeriti.
    \item Il sistema unisce oggetti di base e suggerimenti dai tag, senza duplicati,
    e li aggiunge alla checklist.
\end{enumerate}

\subsubsection*{UC-7: Cerca e filtra elementi}
\textbf{Attore:} Utente.\\
\textbf{Flusso base:} l'utente applica ricerca o filtri (viaggi per
titolo/stato; attività per categoria/stato/giorno; spese per categoria/importo;
tappe per nome/data) e la lista mostra solo i risultati pertinenti.\\
\textbf{Alternativo:} nessun risultato $\rightarrow$ viene mostrato uno stato vuoto
esplicativo.

\subsubsection*{UC-8: Visualizza statistiche}
\textbf{Attore:} Utente.\\
\textbf{Flusso base:} l'utente apre la sezione Statistiche; il sistema calcola
totali, attività completate, spese e avanzamento checklist, e genera i grafici e
le barre di salute del budget.\\
\textbf{Alternativo:} archivio vuoto $\rightarrow$ messaggio ``Nessun dato
disponibile''.

\subsubsection*{UC-9: Visualizza timeline}
\textbf{Attore:} Utente.\\
\textbf{Flusso base:} l'utente apre la scheda Timeline; il sistema mostra le
attività come blocchi colorati per categoria, ordinate per orario, ed evidenzia i
tempi liberi tra un'attività e l'altra.

\subsubsection*{UC-10: Esporta / Importa viaggio}
\textbf{Attore:} Utente.\\
\textbf{Postcondizioni:} Un viaggio è condiviso o ricreato nell'archivio.\\
\textbf{Flusso base (Esporta):} dal dettaglio del viaggio l'utente sceglie
``Esporta''; il sistema genera il JSON e apre il menu di condivisione nativo.\\
\textbf{Flusso base (Importa):} dall'elenco l'utente incolla il JSON; il sistema lo
valida, rigenera gli identificatori, ricollega le relazioni e crea il nuovo viaggio.\\
\textbf{Alternativo:} JSON non valido $\rightarrow$ messaggio d'errore.

\subsubsection*{UC-11: Modifica o elimina elemento}
\textbf{Attore:} Utente.\\
\textbf{Flusso base:} l'utente seleziona ``Modifica'' o ``Elimina''; in caso di
eliminazione il sistema chiede conferma con un \texttt{AlertDialog}, poi aggiorna il
database e la UI.\\
\textbf{Alternativi:} annullamento della rimozione (nessuna modifica) oppure dati
modificati non validi (salvataggio bloccato con segnalazione).

% ============================================================
\section{Conclusioni}
Travel Organizer Offline soddisfa tutte le funzionalità richieste dalla traccia --
viaggi, tappe, attività, checklist e spese -- aggiungendo ricerca, filtri,
statistiche e diverse feature avanzate effettivamente integrate. L'architettura a
livelli, la persistenza locale coerente e i controlli di integrità rendono l'app
robusta, semplice da usare e completamente funzionante offline.

\end{document}
